\begin{solution}{19}
    \begin{subsolution}
        平板受到重力 $\boldsymbol{P}_{\mathrm{C}}$、拉力 $\boldsymbol{Q}_{\mathrm{M}_0}$、铰链对三角形板的作用力 $\boldsymbol{N}_{\mathrm{A}}$ 和 $\boldsymbol{N}_{\mathrm{B}}$，各力及其作用点的坐标分别为：
        \begin{align*}
            \boldsymbol{P}_{\mathrm{C}} &= (0, -mg\sin\varphi, -mg\cos\varphi), \quad (0, 0, h); \\
            \boldsymbol{Q}_{\mathrm{M}_0} &= (0, Q, 0), \quad (x_0, 0, z_0); \\
            \boldsymbol{N}_{\mathrm{A}} &= (N_{\mathrm{Ax}}, N_{\mathrm{Ay}}, N_{\mathrm{Az}}), \quad \left(\frac{b}{2}, 0, 0\right); \\
            \boldsymbol{N}_{\mathrm{B}} &= (N_{\mathrm{B}x}, N_{\mathrm{B}y}, N_{\mathrm{B}z}), \quad \left(-\frac{b}{2}, 0, 0\right)
        \end{align*}
        式中
        \begin{equation*}
            h = \frac{1}{3}\sqrt{a^2 - \frac{b^2}{4}}
        \end{equation*}
        是平板质心到 x 轴的距离。

        平板所受力和（对 O 点的）力矩的平衡方程为
        \begin{align}
            \sum F_x &= N_{\mathrm{Ax}} + N_{\mathrm{Bx}} = 0 \label{eq:1.s19} \\
            \sum F_y &= Q + N_{\mathrm{Ay}} + N_{\mathrm{By}} - mg\sin\varphi = 0 \label{eq:2.s19} \\
            \sum F_z &= N_{\mathrm{A}z} + N_{\mathrm{B}z} - mg\cos\varphi = 0 \label{eq:3.s19} \\
            \sum M_x &= mgh\sin\varphi - Q \cdot z_0 = 0 \label{eq:4.s19} \\
            \sum M_y &= N_{\mathrm{B}z}\frac{b}{2} - N_{\mathrm{A}z}\frac{b}{2} = 0 \label{eq:5.s19} \\
            \sum M_z &= Q \cdot x_0 + N_{\mathrm{Ay}}\frac{b}{2} - N_{\mathrm{By}}\frac{b}{2} = 0 \label{eq:6.s19}
        \end{align}
        联立以上各式解得
        \begin{align*}
            Q &= \frac{mgh\sin\varphi}{z_0}, \\
            N_{\mathrm{Ax}} &= -N_{\mathrm{Bx}}, \\
            N_{Ay} &= \frac{mg\sin\varphi}{2}\left[1 - \frac{h}{b}\left(\frac{b}{z_0} + \frac{2x_0}{z_0}\right)\right], \quad N_{By} = \frac{mg\sin\varphi}{2}\left[1 - \frac{h}{b}\left(\frac{b}{z_0} - \frac{2x_0}{z_0}\right)\right] \\
            N_{\mathrm{A}z} &= N_{\mathrm{B}z} = \frac{1}{2}mg\cos\varphi
        \end{align*}
        即
        \begin{align}
            \boldsymbol{Q}_{\mathrm{M}_0} &= (0, \frac{mgh\sin\varphi}{z_0}, 0), \label{eq:7.s19} \\
            \boldsymbol{N}_{\mathrm{A}} &= \left(N_{\mathrm{Ax}}, \frac{mg\sin\varphi}{2}\left[1 - \frac{h}{b}\left(\frac{b}{z_0} + \frac{2x_0}{z_0}\right)\right], \frac{1}{2}mg\cos\varphi\right), \label{eq:8.s19} \\
            \boldsymbol{N}_{\mathrm{B}} &= \left(-N_{\mathrm{Ax}}, \frac{mg\sin\varphi}{2}\left[1 - \frac{h}{b}\left(\frac{b}{z_0} - \frac{2x_0}{z_0}\right)\right], \frac{1}{2}mg\cos\varphi\right) \label{eq:9.s19}
        \end{align}
    \end{subsolution}
    \begin{subsolution}
        如果希望在 $\mathrm{M}(x,0,z)$ 点的位置从点 $\mathrm{M}_0(x_0,0,z_0)$ 缓慢改变的过程中，可以使铰链支点对板的作用力 $N_{By}$ 保持不变，则需
        \begin{equation}
            N_{By} = \frac{mg\sin\varphi}{2}\left[1 - \frac{h}{b}\left(\frac{b}{z} - \frac{2x}{z}\right)\right] = \text{常量}
            \label{eq:10.s19}
        \end{equation}
        M点移动的起始位置为 $\mathrm{M}_0$，由\eqref{eq:10.s19}式得
        \begin{equation}
            \frac{b}{z} - \frac{2x}{z} = \frac{b}{z_0} - \frac{2x_0}{z_0}
            \label{eq:11.s19}
        \end{equation}
        或
        \begin{equation}
            b - 2x = \left(\frac{b}{z_0} - \frac{2x_0}{z_0}\right)z
            \label{eq:12.s19}
        \end{equation}
        这是过 $\mathrm{A}(\frac{b}{2},0,0)$ 点的直线。

        因此，当力 $\boldsymbol{Q}_{\mathrm{M}}$ 的作用点 M 的位置沿通过 A 点任一条射线（不包含 A 点）在平板上缓慢改变时，铰链支点 B 对板的作用力 $N_{\mathrm{By}}$ 保持不变。同理，当力 $\boldsymbol{Q}_{\mathrm{M}}$ 的作用点 M 沿通过 B 点任一条射线在平板上缓慢改变时，铰链支点 A 对板的作用力 $N_{\mathrm{Ay}}$ 保持不变。
    \end{subsolution}
\end{solution}